\chapter{Parámetros de ejecución de los experimentos}
\label{ap:parametros_experimentos}

En este apéndice se recogen los parámetros de ejecución de los experimentos indicados en el Capítulo~\ref{chap:metodologia_experimental}.

\section{Parámetros de ejecución de las metaheurísticas}

\clearpage
\begingroup
\setlength{\tabcolsep}{5pt}
\renewcommand{\arraystretch}{1.3}

\begin{longtable}[htbp]{|p{3.3cm}|p{2.8cm}|p{5.15cm}|}
    \hline
    \textbf{Parámetro} & \textbf{Valor por defecto}                                         & \textbf{Descripción} \\
    \hline
    \endfirsthead

    \hline
    \textbf{Parámetro} & \textbf{Valor por defecto}                                         & \textbf{Descripción} \\
    \hline
    \endhead

    \texttt{--algoritmo}
                       & \texttt{all}
                       & Metaheurística ejecutada. Valores posibles:
    \begin{itemize}
        \setlength{\itemsep}{0pt}
        \setlength{\topsep}{2pt}
        \item \texttt{age}: algoritmo genético estacionario.
        \item \texttt{de}: evolución diferencial.
        \item \texttt{shade}: variante adaptativa SHADE.
        \item \texttt{all}: ejecuta las tres metaheurísticas.
    \end{itemize}                                                           \\
    \hline

    \texttt{--cec-funcid}
                       & \texttt{-}
                       & Funciones de CEC2017 ejecutadas. Valores posibles:
    \begin{itemize}
        \setlength{\itemsep}{0pt}
        \setlength{\topsep}{2pt}
        \item Lista de identificadores: \texttt{1 4 10}
        \item Lista separada por comas: \texttt{1,4,10}
        \item Suite completa: \texttt{all}
    \end{itemize}                                                                 \\
    \hline

    \texttt{--cec-dim}
                       & \(10\)
                       & Dimensión de las funciones de CEC2017. Valores posibles:
    \begin{itemize}
        \setlength{\itemsep}{0pt}
        \setlength{\topsep}{2pt}
        \item 2
        \item 5
        \item 10
        \item 30
        \item 50
    \end{itemize}                                                                                       \\
    \hline

    \texttt{--seed-start}
                       & \(1\)
                       & Primera semilla si no se usa \texttt{--seeds}.                                            \\
    \hline

    \texttt{--n-seeds}
                       & \(10\)
                       & Semillas consecutivas desde \texttt{--seed-start}.                                        \\
    \hline

    \texttt{--tam-poblacion}
                       & \(50\)
                       & Tamaño de la población.                                                                   \\
    \hline

    \texttt{--max-evals}
                       & \(10000 \cdot D\)
                       & Evaluaciones máximas (\(D\) = dimensión del problema).                                    \\
    \hline

    \texttt{--reinicio}
                       & \texttt{False}
                       & Activa el reinicio elitista si se incluye.                                                \\
    \hline

    \texttt{--reinicio-ratio}
                       & -
                       & Fracción del presupuesto total que debe transcurrir sin una mejora
    relevante del segundo mejor individuo antes de activar el reinicio. Valores evaluados:
    \begin{itemize}
        \setlength{\itemsep}{0pt}
        \setlength{\topsep}{2pt}
        \item \(0.01\)
        \item \(0.03\)
        \item \(0.05\)
        \item \(0.07\)
        \item \(0.10\)
    \end{itemize}                                                                                       \\
    \hline

    \texttt{--outdir}
                       & \makecell[l]{\texttt{results/}                                                            \\\texttt{resultados\_}\\\texttt{benchmark\_mhs}}
                       & Directorio utilizado para almacenar los resultados.                                       \\
    \hline

    \texttt{--seeds}
                       & -
                       & Lista explícita de semillas. Si se utiliza, sustituye a
    \texttt{--seed-start} y \texttt{--n-seeds}. Admite valores separados
    por espacios o comas. Por ejemplo:
    \begin{itemize}
        \setlength{\itemsep}{0pt}
        \setlength{\topsep}{2pt}
        \item \texttt{--seeds 1 2 3}
        \item \texttt{--seeds 1,2,3}
    \end{itemize}                                                                                    \\
    \hline

    \texttt{--sin-metricas}
                       & \texttt{False}
                       & Desactiva el registro de métricas y trayectorias.                                         \\
    \hline

    \texttt{--verbose}
                       & \texttt{False}
                       & Muestra progreso por terminal si se incluye.                                              \\
    \hline
    \caption{Parámetros de \texttt{ejecutar\_metaheuristicas.py}.}
    \label{tab:parametros_mhs_benchmark}
\end{longtable}
\endgroup

\section{Parámetros de ejecución de las técnicas subrogadas}

\clearpage
\begingroup
\setlength{\tabcolsep}{5pt}
\renewcommand{\arraystretch}{1.3}

\begin{longtable}[htbp]{|p{4.2cm}|p{1.8cm}|p{5.25cm}|}
    \hline
    \textbf{Parámetro} & \textbf{Valor por defecto}                                                                   & \textbf{Descripción} \\
    \hline
    \endfirsthead

    \hline
    \textbf{Parámetro} & \textbf{Valor por defecto}                                                                   & \textbf{Descripción} \\
    \hline
    \endhead

    \texttt{--resultados-dir}
                       & -
                       & Directorio con los resultados generados por el programa de ejecución de metaheurísticas.                            \\
    \hline

    \texttt{--algoritmo}
                       & \texttt{all}
                       & Metaheurística que generó los datos. Valores posibles:
    \begin{itemize}
        \setlength{\itemsep}{0pt}
        \setlength{\topsep}{2pt}
        \item \texttt{age}: algoritmo genético estacionario.
        \item \texttt{de}: evolución diferencial.
        \item \texttt{shade}: variante adaptativa SHADE.
        \item \texttt{all}: evalúa las tres metaheurísticas.
    \end{itemize}                                                                                      \\
    \hline

    \texttt{--cec-funcid}
                       & -
                       & Funciones de CEC2017 ejecutadas. Valores posibles:
    \begin{itemize}
        \setlength{\itemsep}{0pt}
        \setlength{\topsep}{2pt}
        \item Lista de identificadores: \texttt{1 4 10}
        \item Lista separada por comas: \texttt{1,4,10}
        \item Suite completa: \texttt{all}
    \end{itemize}                                                                                           \\
    \hline

    \texttt{--modelo}
                       & -
                       & Técnica subrogada a evaluar.                                                                                        \\
    \hline

    \texttt{--seed}
                       & \(42\)
                       & Semilla utilizada por el generador pseudoaleatorio para construir el conjunto de validación.                        \\
    \hline

    \texttt{--max-seeds}
                       & -
                       & Número máximo de trayectorias. Sin límite si se omite.                                                              \\
    \hline

    \texttt{--modelo-params-json}
                       & -
                       & Configuración JSON de hiperparámetros del modelo. Opcional.                                                         \\
    \hline
    \caption[Parámetros de los programas de evaluación \textit{offline}.]{Parámetros de los programas de evaluación \textit{offline} acumulativa y no acumulativa.}
    \label{tab:parametros_runners_offline}
\end{longtable}

\endgroup

\begingroup
\setlength{\tabcolsep}{5pt}
\renewcommand{\arraystretch}{1.3}

\begin{longtable}[htbp]{|p{4.2cm}|p{2.2cm}|p{4.85cm}|}
    \hline
    \textbf{Parámetro} & \textbf{Valor por defecto}                                                    & \textbf{Descripción} \\
    \hline
    \endfirsthead

    \hline
    \textbf{Parámetro} & \textbf{Valor por defecto}                                                    & \textbf{Descripción} \\
    \hline
    \endhead

    \texttt{--trayectoria-dir}
                       & -
                       & Directorio con las trayectorias generadas por las metaheurísticas.                                   \\
    \hline

    \makecell[l]{\texttt{--ajuste-}                                                                                           \\\texttt{resultados-dir}}
                       & \makecell[l]{\scriptsize\texttt{resultados\_}                                                        \\\scriptsize\texttt{benchmark\_}      \\\scriptsize\texttt{surrogates\_}      \\\scriptsize\texttt{offline\_ajuste}}
                       & Subdirectorio donde se almacenan los resultados del ajuste.                                          \\
    \hline

    \texttt{--algoritmo}
                       & \texttt{all}
                       & Metaheurística que generó los datos. Valores posibles:
    \begin{itemize}
        \setlength{\itemsep}{0pt}
        \setlength{\topsep}{2pt}
        \item \texttt{age}: algoritmo genético estacionario.
        \item \texttt{de}: evolución diferencial.
        \item \texttt{shade}: variante adaptativa SHADE.
        \item \texttt{all}: evalúa las tres metaheurísticas.
    \end{itemize}                                                                       \\
    \hline

    \texttt{--cec-funcid}
                       & -
                       & Funciones de CEC2017 utilizadas en la ejecución. Valores posibles:
    \begin{itemize}
        \setlength{\itemsep}{0pt}
        \setlength{\topsep}{2pt}
        \item Identificador individual: \texttt{1} o \texttt{f1}.
        \item Lista separada por comas: \texttt{1,4,10}.
        \item Suite completa: \texttt{all}.
    \end{itemize}                                                                  \\
    \hline

    \texttt{--modelo}
                       & \texttt{all}
                       & Técnica subrogada sometida al ajuste fino. Valores posibles:
    \begin{itemize}
        \setlength{\itemsep}{0pt}
        \setlength{\topsep}{2pt}
        \item \texttt{rbf}
        \item \texttt{svr}
        \item \texttt{mlp}
        \item \texttt{rsm}
        \item \texttt{random\_forest}
        \item \texttt{hgb}
        \item \texttt{lasso}
        \item \texttt{xgboost}
        \item \texttt{all}
    \end{itemize}                                                                                              \\
    \hline

    \texttt{--param-grid-json}
                       & -
                       & Fichero JSON opcional con el espacio de hiperparámetros. Si se omite, se usa
    el grid interno por defecto del modelo.                                                                                   \\
    \hline

    \texttt{--metrica-ajuste}
                       & \texttt{spearman}
                       & Métrica para seleccionar la configuración óptima. Valores posibles:
    \begin{itemize}
        \setlength{\itemsep}{0pt}
        \setlength{\topsep}{2pt}
        \item \texttt{spearman}
        \item \texttt{mae}
        \item \texttt{rmse}
        \item \texttt{nmae}
        \item \texttt{nrmse}
    \end{itemize}                                                                                                  \\
    \hline

    \texttt{--validacion-ratio}
                       & \(0.20\)
                       & Fracción final del bloque para validación interna. Debe cumplir
    \(0 < r < 0.5\).                                                                                                          \\
    \hline

    \texttt{--seed}
                       & \(42\)
                       & Semilla empleada durante el ajuste.                                                                  \\
    \hline

    \makecell[l]{\texttt{--seed-selection-}                                                                                   \\\texttt{random-state}}
                       & \(42\)
                       & Semilla utilizada para seleccionar el subconjunto de trayectorias
    cuando se limita el número de semillas con
    \texttt{--max-seeds}.                                                                                                     \\
    \hline

    \texttt{--max-seeds}
                       & -
                       & Número máximo de trayectorias. Sin límite si se omite. Si se indica, debe ser
    mayor o igual que \(1\) y no superar las trayectorias disponibles.                                                        \\
    \hline
    \caption{Parámetros de \texttt{ajustar\_hiperparametros.py}.}
    \label{tab:parametros_tuning_runner}
\end{longtable}

\endgroup

Los espacios de hiperparámetros
concretos y las configuraciones seleccionadas se presentan junto con sus
resultados en el capítulo siguiente.

\section{Parámetros de ejecución de la hibridación}


\clearpage
\begingroup
\setlength{\tabcolsep}{5pt}
\renewcommand{\arraystretch}{1.3}

\begin{longtable}[htbp]{|p{4.2cm}|p{2.3cm}|p{4.75cm}|}
    \hline
    \textbf{Parámetro} & \textbf{Valor por defecto}                                                     & \textbf{Descripción} \\
    \hline
    \endfirsthead

    \hline
    \textbf{Parámetro} & \textbf{Valor por defecto}                                                     & \textbf{Descripción} \\
    \hline
    \endhead

    \texttt{--algoritmo}
                       & \texttt{age}
                       & Metaheurística \textit{online} ejecutada. Valores posibles:
    \begin{itemize}
        \setlength{\itemsep}{0pt}
        \setlength{\topsep}{2pt}
        \item \texttt{age}
        \item \texttt{de}
        \item \texttt{shade}
        \item \texttt{all}
    \end{itemize}
    También admite listas separadas por espacios o comas.                                                                      \\
    \hline

    \texttt{--cec-funcid}
                       & -
                       & Funciones de CEC2017 ejecutadas. Admite identificadores individuales, listas
    separadas por espacios o comas, y \texttt{all} para la suite completa.                                                     \\
    \hline

    \texttt{--cec-dim}
                       & \(10\)
                       & Dimensión de las funciones CEC2017. Valores posibles: \(2\), \(5\), \(10\),
    \(30\) y \(50\).                                                                                                           \\
    \hline

    \texttt{--seed-start}
                       & \(1\)
                       & Primera semilla si no se proporciona una lista explícita mediante
    \texttt{--seeds}.                                                                                                          \\
    \hline

    \texttt{--n-seeds}
                       & \(10\)
                       & Número de semillas consecutivas generadas a partir de
    \texttt{--seed-start}.                                                                                                     \\
    \hline

    \texttt{--seeds}
                       & -
                       & Lista explícita de semillas. Si se indica, sustituye a
    \texttt{--seed-start} y \texttt{--n-seeds}.                                                                                \\
    \hline

    \texttt{--tam-poblacion}
                       & \(50\)
                       & Tamaño de población utilizado por las metaheurísticas.                                                \\
    \hline

    \texttt{--max-evals}
                       & \(10000 \cdot D\)
                       & Presupuesto máximo de evaluaciones reales.                                                            \\
    \hline

    \texttt{--reinicio}
                       & \texttt{False}
                       & Activa el mecanismo de reinicio elitista en la metaheurística \textit{online}.                        \\
    \hline

    \texttt{--reinicio-ratio}
                       & -
                       & Fracción del presupuesto que debe transcurrir sin mejora antes de permitir
    el reinicio. Solo puede utilizarse junto con \texttt{--reinicio}.                                                          \\
    \hline

    \texttt{--modelo}
                       & Modelo seleccionado
                       & Modelo subrogado empleado por el filtro \textit{online}.                                              \\
    \hline

    \makecell[l]{\texttt{--modelo-params-}                                                                                     \\\texttt{json}}
                       & Config. del modelo seleccionado
                       & Fichero JSON con los hiperparámetros del modelo subrogado.                                            \\
    \hline

    \texttt{--modelo-prob}
                       & \(0.50\)
                       & Probabilidad \(p\) de aplicar el filtro subrogado a un candidato. Debe estar
    en \([0,1]\). En las pruebas se consideran \(0.00\), \(0.25\), \(0.50\)
    y \(0.75\).                                                                                                                \\
    \hline

    \texttt{--warmup-ratio}
                       & \(0.20\)
                       & Proporción inicial del presupuesto durante la cual el subrogado permanece
    inactivo y todas las evaluaciones aceptadas se realizan con la función real.                                               \\
    \hline

    \texttt{--window-ratio}
                       & \(0.20\)
                       & Proporción del presupuesto utilizada para definir la ventana reciente de
    entrenamiento del subrogado.                                                                                               \\
    \hline

    \texttt{--retrain-ratio}
                       & \(0.25\)
                       & Fracción de la ventana que debe renovarse antes de reentrenar el modelo.
    En las pruebas se consideran \(0.25\), \(0.50\) y \(0.75\).                                                                \\
    \hline

    \makecell[l]{\texttt{--cooldown-}                                                                                          \\\texttt{reinicio-evals}}
                       & \(0\)
                       & Número de evaluaciones reales de enfriamiento tras cada reinicio. El valor
    \(0\) desactiva este mecanismo. En las pruebas se consideran \(0\), \(50\),
    \(100\) y \(500\).                                                                                                         \\
    \hline

    \makecell[l]{\texttt{--guardar-}                                                                                           \\\texttt{decisiones-subrogado}}
                       & \texttt{False}
                       & Guarda un fichero con las decisiones del filtro, incluyendo la predicción,
    la referencia real y el margen entre ambas.                                                                                \\
    \hline

    \texttt{--outdir}
                       & \makecell[l]{\texttt{results/}                                                                        \\\texttt{cec/}                           \\\texttt{cec2017\_d10}                  \\\texttt{\_tam50}                       \\\texttt{\_online}}
                       & Directorio donde se almacenan los resultados de la ejecución \textit{online}.                         \\
    \hline

    \texttt{--verbose}
                       & \texttt{False}
                       & Muestra el progreso de las ejecuciones por terminal.                                                  \\
    \hline
    \caption{Parámetros de \texttt{ejecutar\_metaheuristicas\_online.py}.}
    \label{tab:parametros_runner_online}
\end{longtable}

\endgroup








